% Audit de veracite des citations de rapport de jury - 2026-09-20
% Bloc 4 : citation jury 2024 supprimee (introuvable dans le corpus).
% Le remplacement envisage (RJ2023 l.459) aurait fait doublon exact avec la
% citation jury 2023 deja presente, qui mentionne deja la marge de phase.
\expandafter\ifx\csname CAPETINCLUDEMODE\endcsname\relax
\documentclass[11pt,a4paper]{article}
% preamble.tex — packages et commandes communes a toutes les fiches CAPET SII IE
\input{glyphtounicode}
\pdfgentounicode=1
\usepackage{cmap}
\usepackage[T1]{fontenc}
\usepackage[utf8]{inputenc}
\usepackage[french]{babel}
% babel[french] active a \begin{document} des raccourcis de ponctuation (:;!?)
% qui cassent la syntaxe to[...] de circuitikz (erreur "+ or - expected") ;
% on les desactive, ce qui ne change que l'espacement fin avant ces signes.
\AtBeginDocument{\shorthandoff{:;!?}}
\usepackage[a4paper, top=1.6cm, bottom=1.6cm, left=1.4cm, right=1.4cm]{geometry}
\setlength{\headheight}{14.5pt}
\usepackage{amsmath, amssymb, amsfonts}
\usepackage{siunitx}
\sisetup{locale=FR, output-decimal-marker={,}, inter-unit-product=\ensuremath{{}\cdot{}},
         group-minimum-digits=4}
% Unites non fournies par siunitx, declarees ici une fois pour toutes (fiches + exercices).
% Les gardes \ifdefined evitent tout conflit avec une redeclaration locale restee dans un
% fichier : ne pas retirer les gardes, ne pas redeclarer ces unites dans les fichiers.
\DeclareSIUnit{\jour}{j}
\ifdefined\var\else\DeclareSIUnit{\var}{var}\fi   % puissance reactive
\ifdefined\vad\else\DeclareSIUnit{\vad}{VAD}\fi   % puissance deformante
\ifdefined\tr\else\DeclareSIUnit{\tr}{tr}\fi      % tour (tr/min)
\usepackage{array, booktabs, multirow, tabularx}
\usepackage{xcolor}
\usepackage[most, breakable]{tcolorbox}
\usepackage{multicol}
\usepackage{enumitem}
\usepackage{mathtools}
\usepackage{lmodern}
% Refonte visuelle 2026 (direction "moderne coloré") : Lato pour le texte,
% maths sans-serif assorties via newtxsf (charge apres amsmath/amssymb).
\usepackage[default]{lato}
\usepackage{newtxsf}
\usepackage{microtype}
\usepackage{fancyhdr}
\usepackage{tikz}
\usetikzlibrary{positioning,arrows.meta}
\usepackage{circuitikz}
\usepackage{lastpage}

% ── Couleurs par domaine (palette modernisée 2026) ────────────────────
\definecolor{colElec}{RGB}{29, 78, 216}        % Électrotechnique — bleu vif
\definecolor{colElecLight}{RGB}{219, 234, 254} % fond bleu clair
\definecolor{colAuto}{RGB}{21, 128, 61}        % Automatique — vert
\definecolor{colMeca}{RGB}{194, 65, 12}        % Mécanique — orange
\definecolor{colTherm}{RGB}{185, 28, 28}       % Thermique — rouge
\definecolor{colQT}{RGB}{232, 244, 233}        % (hérité, plus utilisé par questiontype)
\definecolor{colJury}{RGB}{255, 248, 225}      % fond jury (crème)
\definecolor{colJuryFrame}{RGB}{202, 138, 4}   % bordure/bandeau jury (ambre)
\definecolor{colGrayLight}{RGB}{245, 245, 247} % fond gris neutre
\definecolor{colSeparator}{RGB}{180, 180, 180} % gris séparateur
\definecolor{colExo}{RGB}{8, 145, 178}          % Banque d'exercices — cyan (distinct des fiches)
% ── Encre et gris de la refonte ───────────────────────────────────────
\definecolor{colInk}{RGB}{23, 32, 42}          % texte courant (quasi-noir)
\definecolor{colSlate}{RGB}{83, 98, 118}       % gris ardoise (méta, en-têtes)
\definecolor{colRule}{RGB}{203, 213, 225}      % filets clairs
\definecolor{colJuryBg}{RGB}{255, 248, 225}    % fond crème jury
\definecolor{colPresence}{RGB}{83, 98, 118}    % bloc 5 (présence dans les sujets)
\AtBeginDocument{\color{colInk}}

% ── Style encadré exercice + corrigé (banque d'exercices) ─────────────
\tcbset{
  exoBox/.style={
    breakable,
    enhanced,
    colback=colExo!8!white,
    colframe=colExo,
    colbacktitle=colExo,
    coltitle=white,
    fonttitle=\bfseries\small,
    left=4pt, right=4pt, top=3pt, bottom=3pt,
    before skip=5pt, after skip=5pt,
  }
}

% ── Bascule enonce / corrige (banque d'exercices) ─────────────────────
% \CAPETCORRIGE est defini en ligne de commande (pdflatex "\def\CAPETCORRIGE{1}\input{...}")
% pour compiler la variante corriges d'un meme fichier source ; par defaut (non defini),
% la compilation produit la variante enonces. \CAPETSHOW{enonce}{corrige} choisit la branche.
%
% TROIS MODES (refonte 2026-09-20, manuel de revision) :
%   - aucun drapeau   -> imprime l'enonce            (variante enonces, INCHANGE)
%   - \CAPETCORRIGE   -> imprime le corrige          (variante corriges, INCHANGE)
%   - \CAPETMANUEL    -> imprime l'enonce ET met le corrige de cote pour l'annexe
%
% Le mode manuel stocke chaque corrige dans une macro (\csgdef), sans ecriture
% fichier ni re-tokenisation : le contenu (circuitikz, accents, %) traverse intact.
% L'annexe les rejoue ensuite dans l'ordre via \CAPETANNEXE.
\usepackage{etoolbox}

\newcounter{CAPETexo}          % numero continu d'exercice dans le manuel
\setcounter{CAPETexo}{0}
\newcounter{CAPETannexe}       % compteur de relecture en annexe

% Titre courant de la notion en cours (pose par \CAPETNOTION), sert d'intertitre annexe
\newcommand{\CAPETnotioncode}{}
\newcommand{\CAPETnotiontitre}{}

% \CAPETNOTION{code}{titre} : ouvre une notion dans le manuel (memorise le libelle,
% et marque le rang du premier exercice de la notion pour l'intertitre d'annexe).
\newcommand{\CAPETNOTION}[2]{%
  \renewcommand{\CAPETnotioncode}{#1}%
  \renewcommand{\CAPETnotiontitre}{#2}%
  \csgdef{CAPETnotionstart\the\numexpr\value{CAPETexo}+1\relax}{#1 -- #2}%
}

% ── \degree robuste ───────────────────────────────────────────────────
% Certains fichiers d'exercices utilisent \degree. Ni le preambule ni siunitx
% ne le definissent, et \providecommand local ne suffit pas : sous hyperref,
% \degree existe deja comme commande factice (\ERROR), donc \providecommand
% s'abstient et l'appel echoue a la compilation du manuel (erreur invisible
% en compilation isolee). On le (re)definit donc inconditionnellement ici,
% apres le chargement des packages, pour tout le corpus.
\AtBeginDocument{\providecommand{\degree}{}\renewcommand{\degree}{\ensuremath{^\circ}}}

% Libelle decoratif : les fichiers d'exercices utilisent aussi \CAPETSHOW dans
% leur \fancyhead et leur bloc de titre, pour afficher "(enonces)" ou "(corriges)".
% Ces appels-la ne sont PAS des exercices : en mode manuel ils ne doivent ni
% consommer un numero, ni alimenter l'annexe. On les reconnait a leur premier
% argument, strictement egal au libelle ci-dessous.
\newcommand{\CAPETLIBELLEENONCES}{(\'enonc\'es)}

\ifdefined\CAPETMANUEL
  % --- Mode manuel : enonce en place, corrige differe vers l'annexe ---
  \newcommand{\CAPETSHOW}[2]{%
    \def\CAPETargun{#1}%
    \ifx\CAPETargun\CAPETLIBELLEENONCES
      #1%            appel decoratif : rendu tel quel, hors numerotation
    \else
      \CAPETSHOWEXO{#1}{#2}%
    \fi
  }
  \newcommand{\CAPETSHOWEXO}[2]{%
    \stepcounter{CAPETexo}%
    \csgdef{CAPETcorrige\theCAPETexo}{#2}%
    \edef\CAPETtmpnum{\theCAPETexo}%
    #1%
    \par\vspace{2pt}%
    {\footnotesize\itshape\color{colExo}%
     \hyperref[capet:cor:\CAPETtmpnum]{Corrige de l'exercice \CAPETtmpnum\ $\rightarrow$ p.~\pageref{capet:cor:\CAPETtmpnum}}%
     \label{capet:eno:\CAPETtmpnum}}%
  }
  % \CAPETANNEXE : rejoue tous les corriges captures, dans l'ordre, avec intertitres
  \newcommand{\CAPETANNEXE}{%
    \setcounter{CAPETannexe}{0}%
    \loop\ifnum\value{CAPETannexe}<\value{CAPETexo}%
      \stepcounter{CAPETannexe}%
      \edef\CAPETcurnum{\theCAPETannexe}%
      \ifcsname CAPETnotionstart\CAPETcurnum\endcsname
        \par\medskip
        {\normalsize\bfseries\color{colExo}\csuse{CAPETnotionstart\CAPETcurnum}}%
        \par\nobreak\smallskip
      \fi
      \begin{tcolorbox}[exoBox, breakable,
        title={Corrige de l'exercice \CAPETcurnum\
               \hfill\normalfont\footnotesize
               \hyperref[capet:eno:\CAPETcurnum]{enonce p.~\pageref{capet:eno:\CAPETcurnum}}}]%
      \label{capet:cor:\CAPETcurnum}%
      \csuse{CAPETcorrige\CAPETcurnum}%
      \end{tcolorbox}%
    \repeat
  }
\else
  % --- Modes historiques, comportement strictement inchange ---
  \newcommand{\CAPETSHOW}[2]{\ifdefined\CAPETCORRIGE #2\else #1\fi}
  \newcommand{\CAPETANNEXE}{}
\fi

% ── Difficulté 1 à 3 (étoiles pleines, amssymb) ───────────────────────
\newcommand{\difficulte}[1]{%
  \ifcase#1\relax
  \or $\bigstar$\;niveau 1/3%
  \or $\bigstar\bigstar$\;niveau 2/3%
  \or $\bigstar\bigstar\bigstar$\;niveau 3/3%
  \fi
}

% ── Style encadré extrait de document technique (banque d'exercices) ──
\tcbset{
  docTech/.style={
    breakable,
    enhanced,
    colback=colGrayLight,
    colframe=colSeparator,
    fonttitle=\bfseries\small\itshape,
    coltitle=black,
    colbacktitle=colGrayLight,
    left=4pt, right=4pt, top=3pt, bottom=3pt,
    before skip=4pt, after skip=4pt,
  }
}

% ── Style tcolorbox de section principale ────────────────────────────
% Refonte 2026 : bandeau titre coloré arrondi + fond très pâle, sans cadre.
\tcbset{
  ficheSection/.style={
    breakable, enhanced jigsaw,
    frame hidden, boxrule=0pt, arc=1.8mm,
    colback=#1!4!white,
    colbacktitle=#1,
    coltitle=white,
    fonttitle=\sffamily\bfseries,
    left=7pt, right=7pt, top=5pt, bottom=5pt,
    toptitle=3pt, bottomtitle=3pt,
    before skip=10pt, after skip=8pt,
  }
}

% ── Style encadré question type ──────────────────────────────────────
% Carte blanche, bandeau sombre, liseré très clair.
\tcbset{
  questiontype/.style={
    breakable, enhanced jigsaw,
    colback=white,
    frame hidden, boxrule=0pt, arc=1.8mm,
    colbacktitle=colElec!78!black,
    coltitle=white,
    fonttitle=\sffamily\bfseries\small,
    left=7pt, right=7pt, top=4pt, bottom=4pt,
    toptitle=2.5pt, bottomtitle=2.5pt,
    before skip=8pt, after skip=8pt,
    borderline={0.6pt}{0pt}{colElec!12!white},
  }
}

% ── Style encadré jury ───────────────────────────────────────────────
% Bandeau ambre, fond crème, sans cadre.
\tcbset{
  jury/.style={
    breakable, enhanced jigsaw,
    colback=colJuryBg,
    frame hidden, boxrule=0pt, arc=1.8mm,
    colbacktitle=colJuryFrame,
    coltitle=white,
    fonttitle=\sffamily\bfseries,
    left=7pt, right=7pt, top=4pt, bottom=4pt,
    toptitle=3pt, bottomtitle=3pt,
    before skip=10pt, after skip=8pt,
  }
}

% ── En-tête / pied de page communs (refonte 2026) ─────────────────────
% #1 code fiche, #2 titre court, #3 couleur domaine
\newcommand{\ficheHeaders}[3]{%
  \fancyhead[L]{\small\sffamily{\bfseries\textcolor{#3}{#1}}\hspace{0.7em}\textcolor{colSlate}{#2}}%
  \fancyhead[R]{\small\sffamily\textcolor{colSlate}{CAPET SII IE\, --- \,Session 2026}}%
  \fancyfoot[C]{\small\sffamily\textcolor{colSlate}{\thepage\,/\,\pageref{LastPage}}}%
  \fancyhead[C]{}\fancyfoot[L]{}\fancyfoot[R]{}%
}

% ── Bloc en-tête de fiche (BLOC 0, refonte 2026) ──────────────────────
% #1 code  #2 titre  #3 couleur domaine  #4 domaine  #5 ligne présence  #6 sous-thèmes
\newcommand{\ficheEntete}[6]{%
  \begin{tcolorbox}[enhanced, breakable,
    frame hidden, boxrule=0pt, arc=2.5mm,
    interior style={left color=#3, right color=#3!75!black},
    left=9pt, right=9pt, top=7pt, bottom=7pt,
    before skip=0pt, after skip=6pt]
    {\color{white}\sffamily\bfseries\LARGE #1\quad #2}\\[4pt]
    {\color{#3!22!white}\sffamily\small \MakeUppercase{#4}\quad\textbar\quad #5}\\[1pt]
    {\color{#3!35!white}\sffamily\footnotesize Sous-thèmes : #6}
  \end{tcolorbox}%
}

% ── Mise en page ─────────────────────────────────────────────────────
\pagestyle{fancy}
\fancyhf{}
\renewcommand{\headrulewidth}{0.5pt}
\renewcommand{\footrulewidth}{0pt}
\AtBeginDocument{%
  \renewcommand{\headrule}{{\color{colRule}\hrule width\headwidth height 0.5pt}}%
}

\setlist[itemize]{leftmargin=1.2em, itemsep=1pt, topsep=2pt, parsep=0pt}
\setlist[enumerate]{leftmargin=1.2em, itemsep=3pt, topsep=2pt, parsep=0pt}

\setlength{\parskip}{2pt}
\setlength{\columnsep}{10pt}

\def\CAPETBEGINFILE{\begin{document}}
\def\CAPETENDFILE{\end{document}}
\else
\def\CAPETBEGINFILE{}
\def\CAPETENDFILE{}
\fi

\ficheHeaders{T27}{Correcteurs et stabilité}{colAuto}

\CAPETBEGINFILE

% -----------------------------------------------------------------------
%  BLOC EN-TETE
% -----------------------------------------------------------------------

\ficheEntete{T27}{Systèmes bouclés, correcteurs, stabilité}{colAuto}%
  {Automatique - Asservissement et régulation}%
  {\textbf{Présence :} \textbf{9/13} sujets (2013-2026) - thème \emph{très fréquent}}%
  {FTBO / FTBF - marges de phase et de gain - critère de Bode - correcteurs P, PI, PID - réglage graphique - correcteur numérique (récurrence)}

% -----------------------------------------------------------------------
%  1. RAPPELS DE COURS
% -----------------------------------------------------------------------

\begin{tcolorbox}[
  ficheSection=colAuto,
  title={1. Rappels de cours},
]

% -- 1.1 Structure d'un asservissement -------------------------------
\textbf{1.1 Structure d'un asservissement à retour unitaire}

\smallskip
\begin{center}
\begin{tikzpicture}[>=Stealth, font=\small]
  \node (E) at (0,0) {$E(p)$};
  \node[draw, circle, minimum size=6mm, inner sep=0pt] (sum) at (1.7,0) {};
  \draw (sum.north) -- (sum.south);
  \draw (sum.east) -- (sum.west);
  \node[font=\tiny] at ($(sum.center)+(-0.18,0.18)$) {$+$};
  \node[font=\tiny] at ($(sum.center)+(-0.18,-0.18)$) {$-$};
  \node[draw, minimum width=13mm, minimum height=8mm] (C) at (3.6,0) {$C(p)$};
  \node[draw, minimum width=17mm, minimum height=8mm] (Hproc) at (6.0,0) {$H_{proc}(p)$};
  \node (S) at (8.0,0) {$S(p)$};
  \node[draw, minimum width=17mm, minimum height=8mm] (Hcapt) at (6.0,-1.4) {$H_{capt}(p)$};

  \draw[->] (E) -- (sum);
  \draw[->] (sum) -- node[above]{$\varepsilon(p)$} (C);
  \draw[->] (C) -- (Hproc);
  \draw[->] (Hproc) -- (S);
  \fill (7.2,0) circle (0.9pt);
  \draw[->] (7.2,0) |- (Hcapt.east);
  \draw[->] (Hcapt.west) -| (sum.south);
\end{tikzpicture}
\end{center}

\smallskip
\renewcommand{\arraystretch}{1.2}
\begin{center}
\begin{tabular}{ll}
\toprule
\textbf{Grandeur} & \textbf{Expression} \\
\midrule
FTBO & $H_{BO}(p) = C(p)\,H_{proc}(p)\,H_{capt}(p)$ \\
FTBF (retour unitaire) & $H_{BF}(p) = \dfrac{H_{BO}(p)}{1 + H_{BO}(p)}$ \\
Erreur en boucle fermée & $\varepsilon(p) = \dfrac{E(p)}{1 + H_{BO}(p)}$ \\
\bottomrule
\end{tabular}
\end{center}

\medskip\hrule\vspace{4pt}

% -- 1.2 Performances ------------------------------------------------
\textbf{1.2 Performances en régime permanent}

\begin{itemize}
  \item \textbf{Erreur statique (échelon)} : $\varepsilon_s = \lim_{p\to 0} p\,\varepsilon(p) = \dfrac{E_0}{1 + H_{BO}(0)}$
  \item Si la FTBO contient au moins un intégrateur : $H_{BO}(0) \to \infty$ donc $\varepsilon_s = 0$
  \item \textbf{Dépassement} : $D\% = 100\,\exp\!\left(\dfrac{-m\pi}{\sqrt{1-m^2}}\right)$ \quad ($m < 1$)
  \item \textbf{Temps de réponse à 5\,\%} : $t_{R5\%} \approx 3\tau$ (1\textsuperscript{er} ordre)
\end{itemize}

\medskip\hrule\vspace{4pt}

% -- 1.3 Stabilite - marges (Bode) -----------------------------------
\textbf{1.3 Stabilité - critère des marges (diagramme de Bode de la BO)}

\begin{center}
\renewcommand{\arraystretch}{1.3}
\begin{tabular}{lll}
\toprule
\textbf{Marge} & \textbf{Définition} & \textbf{Lecture graphique} \\
\midrule
Marge de phase $M_\varphi$ & $M_\varphi = 180^{\circ} + \arg(H_{BO}(j\omega_c))$ & Phase à $\omega_c$ (gain = 0 dB) \\
Marge de gain $M_G$ & $M_G = -|H_{BO}(j\omega_{-180^{\circ}})|_{dB}$ & Gain à $\omega_{-180^{\circ}}$ (phase = $-180^{\circ}$) \\
\bottomrule
\end{tabular}
\end{center}

\textbf{Cahiers des charges rencontrés :} $M_\varphi \geq 45^{\circ}$ (2023, 2026) ; $M_\varphi \geq 90^{\circ}$ (2024) ; $M_G \geq 10$ dB (2026).

\textbf{Condition de stabilité :} $M_\varphi > 0^{\circ}$ \textit{et} $M_G > 0$ dB : système stable en boucle fermée.

\medskip\hrule\vspace{4pt}

% -- 1.4 Correcteurs -------------------------------------------------
\textbf{1.4 Correcteurs courants}

\medskip
\begin{center}
\renewcommand{\arraystretch}{1.35}
\begin{tabular}{llp{6cm}}
\toprule
\textbf{Type} & \textbf{Transmittance} & \textbf{Effet} \\
\midrule
P & $C(p) = K_C$ & Gain pur : améliore la précision, réduit $M_\varphi$ \\
PI & $K_P\!\left(1 + \dfrac{1}{T_i p}\right)$ & Annule $\varepsilon_s$ (intégrateur) ; ajoute $-90^{\circ}$ de phase \\
PID & $K_P\!\left(1 + \dfrac{1}{T_i p} + T_d p\right)$ & P (rapidité) + I (précision) + D (stabilité) \\
\bottomrule
\end{tabular}
\end{center}

\medskip\hrule\vspace{4pt}

% -- 1.5 Reglage du correcteur par Bode ------------------------------
\textbf{1.5 Réglage du correcteur P par la méthode de Bode}

\textbf{Procédure :}
\begin{enumerate}
  \item Tracer le Bode de $H_{BO}(j\omega)$ pour $K_C = 1$.
  \item Identifier $\omega_B$ tel que $\arg(H_{BO}(j\omega_B)) = -180^{\circ} + M_{\varphi\,\text{cible}}$.
  \item Le gain à $\omega_B$ vaut $-G_{dB}$ ; pour obtenir 0 dB à $\omega_B$ :
  \[
    K_{C\,\lim} = 10^{G_{dB}/20}
  \]
\end{enumerate}

\textit{Exemple (corrigé 2023) :} $M_\varphi = 45^{\circ}$, $\omega_B = 7\,000$ rad$\cdot$s$^{-1}$, gain $= -59{,}4$ dB $\Rightarrow K_{PV} = 10^{59{,}4/20} = 933$.

\medskip\hrule\vspace{4pt}

% -- 1.6 Correcteur numerique -----------------------------------------
\textbf{1.6 Correcteur PI numérique - relation de récurrence}

Pour $C(p) = K_{co}\dfrac{1 + T_{co}\,p}{T_{co}\,p}$, avec période d'échantillonnage $T_e$ :
\[
  u_k = u_{k-1} + a_0\,\varepsilon_k + a_1\,\varepsilon_{k-1}
\]
\[
  a_0 = K_{co}\!\left(1 + \frac{T_e}{T_{co}}\right) \qquad a_1 = -K_{co}
\]

\textit{Source : E2 Sujet 2026 Q.B19 - asservissement de température d'une pile à combustible.}

\end{tcolorbox}

% -----------------------------------------------------------------------
%  2. FORMULAIRE SYNTHETIQUE
% -----------------------------------------------------------------------

\begin{tcolorbox}[
  ficheSection=colAuto,
  title={2. Formulaire synthétique},
]
\small
\begin{multicols}{2}
\renewcommand{\arraystretch}{1.45}
\begin{tabular}{@{}p{4.8cm}ll@{}}
\toprule
\textbf{Formule} & \textbf{Usage} & \textbf{U.} \\
\midrule
\multicolumn{3}{@{}l}{\footnotesize\textit{Boucle fermée}} \\
$H_{BF} = H_{BO}/(1+H_{BO})$ & FTBF & - \\
$\varepsilon_s = E_0/(1+H_{BO}(0))$ & Erreur statique & - \\
$D\% = 100\,e^{-m\pi/\sqrt{1-m^2}}$ & Dépassement & \% \\
\midrule
\multicolumn{3}{@{}l}{\footnotesize\textit{Marges (Bode)}} \\
$M_\varphi = 180^{\circ} + \arg H_{BO}(j\omega_c)$ & Phase & $^{\circ}$ \\
$M_G = -|H_{BO}(j\omega_{-180^{\circ}})|_{dB}$ & Gain & dB \\
\bottomrule
\end{tabular}

\columnbreak

\renewcommand{\arraystretch}{1.45}
\begin{tabular}{@{}p{4.8cm}ll@{}}
\toprule
\textbf{Formule} & \textbf{Usage} & \textbf{U.} \\
\midrule
\multicolumn{3}{@{}l}{\footnotesize\textit{Correcteurs}} \\
$C_P(p) = K_C$ & Gain pur & - \\
$C_{PI}(p) = K_P(1+1/(T_i p))$ & PI & - \\
$K_{C\,\lim} = 10^{G_{dB}/20}$ & Réglage P & - \\
\midrule
\multicolumn{3}{@{}l}{\footnotesize\textit{Numérique}} \\
$u_k = u_{k-1} + a_0\varepsilon_k + a_1\varepsilon_{k-1}$ & Récurrence PI & - \\
$a_0 = K_{co}(1+T_e/T_{co})$ & Coeff. & - \\
$a_1 = -K_{co}$ & Coeff. & - \\
\bottomrule
\end{tabular}
\end{multicols}
\normalsize
\end{tcolorbox}

% -----------------------------------------------------------------------
%  3. QUESTIONS TYPES
% -----------------------------------------------------------------------

\begin{tcolorbox}[
  ficheSection=colAuto,
  title={3. Questions types - extraites des sujets réels},
]

% -- QT1 -------------------------------------------------------------
\begin{tcolorbox}[questiontype,
  title={QT1 - Lecture graphique des marges \hfill\normalfont\footnotesize E1 2026 (C5)}
]
\textbf{Énoncé.} À partir du diagramme de Bode de la FTBO fourni en document réponse, déterminer graphiquement la marge de phase et la marge de gain. Conclure sur le respect du cahier des charges ($M_\varphi > 45^{\circ}$, $M_G > 10$ dB).

\tcbline

\textbf{Corrigé.}

\textbf{Étape 1 — Marge de phase.}
\begin{enumerate}
  \item Trouver $\omega_c$ tel que $|H_{BO}(j\omega_c)|_{dB} = 0$ dB sur la courbe de gain.
  \item Lire la phase à $\omega_c$ : soit $\varphi(\omega_c)$.
  \item $M_\varphi = 180^{\circ} + \varphi(\omega_c)$ \quad [attention au signe : la phase est négative]
\end{enumerate}

\textbf{Étape 2 — Marge de gain.}
\begin{enumerate}
  \item Trouver $\omega_{-180^{\circ}}$ tel que $\varphi(\omega_{-180^{\circ}}) = -180^{\circ}$.
  \item Lire le gain à $\omega_{-180^{\circ}}$ : soit $G_{dB}$.
  \item $M_G = -G_{dB}$ (en dB, $G_{dB}$ est négatif pour un système stable)
\end{enumerate}

\textbf{Conclusion.} Comparer aux valeurs du CdC : si $M_\varphi > 45^{\circ}$ \textbf{et} $M_G > 10$ dB, le cahier des charges est respecté ; sinon, le correcteur doit être révisé.
\end{tcolorbox}

\vspace{4pt}

% -- QT2 -------------------------------------------------------------
\begin{tcolorbox}[questiontype,
  title={QT2 - Détermination graphique de $K_{C\,\lim}$ \hfill\normalfont\footnotesize E1 2026 (C6)}
]
\textbf{Énoncé.}
Sur le diagramme de Bode de la FTBO avec $K_C = 1$, déterminer $K_{C\,\lim}$ pour se placer à la limite du critère de marge de phase ($M_\varphi = 45^{\circ}$).

\tcbline

\textbf{Corrigé.}

\textbf{Étape 1 — Pulsation cible.} Identifier $\omega_B$ tel que $\varphi(\omega_B) = -180^{\circ} + 45^{\circ} = -135^{\circ}$.

\textbf{Étape 2 — Lecture du gain.} Lire le gain à $\omega_B$ : $|H_{BO}(j\omega_B)|_{dB} = -G_{dB}$ (valeur négative).

\textbf{Étape 3 — Réglage.} Pour que le gain soit exactement 0 dB à $\omega_B$, il faut décaler la courbe de $+G_{dB}$ :
\[
  K_{C\,\lim} = 10^{G_{dB}/20}
\]

\textit{Rappel :} multiplier $K_C$ décale la courbe de gain de $20\log_{10} K_C$ sans modifier la courbe de phase.

\textbf{Conclusion.} Avec $K_C = K_{C\,\lim}$, le système fonctionne exactement à la marge de phase limite imposée par le cahier des charges : tout réglage plus agressif ($K_C$ plus grand) ferait chuter $M_\varphi$ en dessous du seuil requis.
\end{tcolorbox}

\vspace{4pt}

% -- QT3 -------------------------------------------------------------
\begin{tcolorbox}[questiontype,
  title={QT3 - Relation de récurrence du correcteur numérique \hfill\normalfont\footnotesize E2 2026 (B18-B19)}
]
\textbf{Énoncé.}
Un correcteur PI de température d'une pile à combustible a pour transmittance $C(p) = K_{co}\dfrac{1 + T_{co}\,p}{T_{co}\,p}$. La période d'échantillonnage est $T_e$. Établir l'équation différentielle liant $u(t)$ et $\varepsilon(t)$, puis la relation de récurrence.

\tcbline

\textbf{Corrigé.}

\textbf{Étape 1 — Équation différentielle.} À partir de $U(p) = K_{co}\dfrac{1+T_{co}\,p}{T_{co}\,p}\,\mathcal{E}(p)$ :
\[
  T_{co}\,\frac{du}{dt} = K_{co}\,T_{co}\,\frac{d\varepsilon}{dt} + K_{co}\,\varepsilon(t)
\]

\textbf{Étape 2 — Discrétisation} (approximation d'Euler arrière : $\dot{f}(kT_e) \approx (f_k - f_{k-1})/T_e$) :
\[
  u_k = u_{k-1} + a_0\,\varepsilon_k + a_1\,\varepsilon_{k-1}
  \qquad
  a_0 = K_{co}\!\left(1 + \frac{T_e}{T_{co}}\right) \quad a_1 = -K_{co}
\]

\textbf{Conclusion.} Cette relation de récurrence est directement implantable dans le calculateur du régulateur de température : elle respecte la contrainte temps réel imposée par la période d'échantillonnage $T_e$ du sujet.
\end{tcolorbox}

\end{tcolorbox}

% -----------------------------------------------------------------------
%  4. ATTENTES DU JURY
% -----------------------------------------------------------------------

\begin{tcolorbox}[
  jury,
  title={4. Attentes du jury - Erreurs récurrentes et points de vigilance},
]

\begin{enumerate}[label=\textbf{\arabic*.}, leftmargin=*, itemsep=5pt]

  \item \textbf{Lire les marges sur le bon diagramme.} La marge de phase se lit à la pulsation de coupure du \emph{gain} (0 dB), \textbf{pas} à la pulsation de déphasage. La confusion est fréquente.

  \item \textbf{Conclure vis-à-vis du cahier des charges.} Donner une valeur numérique sans la comparer à la valeur cible est insuffisant : écrire « $M_\varphi = 52^{\circ} > 45^{\circ}$ : CDC respecté ».

  \item \textbf{Effet du correcteur P sur la phase.} Augmenter $K_C$ décale la courbe de gain vers le haut ($+20\log K_C$) \textbf{sans toucher la phase} : $\omega_c$ se déplace, $M_\varphi$ change.

  \item \textbf{Correcteur PI : l'intégrateur ajoute $-90^{\circ}$.} L'action intégrale garantit $\varepsilon_s = 0$ pour un échelon mais dégrade la marge de phase. Toujours vérifier la stabilité après introduction du PI.

  \item \textbf{Relation de récurrence : signe des coefficients.} $a_1 = -K_{co}$ est négatif. Erreur fréquente : utiliser $\dot{g}(kT_e) \approx (g_{k+1}-g_k)/T_e$ (Euler avant) au lieu de $(g_k-g_{k-1})/T_e$ (Euler arrière).

  \item \textbf{Thème systématiquement présent depuis 2022.} Ne pas faire l'impasse : les parties asservissement sont fortement pondérées et très pénalisées si absentes.

\end{enumerate}

\medskip
\begin{tcolorbox}[
  colback=colJuryFrame!12!colJury,
  colframe=colJuryFrame,
  left=4pt, right=4pt, top=2pt, bottom=2pt,
  fontupper=\small\itshape
]
Rapport jury 2023 : «~Trop peu de candidats l'ont traitée de manière satisfaisante. Le jury invite les candidats à revoir les concepts de base de ce domaine comme les notions de boucle ouverte, de correcteur ou encore de marge de phase.~»
\end{tcolorbox}

\end{tcolorbox}

% -----------------------------------------------------------------------
%  PIED - PRESENCE DANS LES SUJETS
% -----------------------------------------------------------------------

\begin{tcolorbox}[
  ficheSection=colPresence,
  title={Présence dans les sujets analysés (2013-2026)},
]
\begin{center}
\footnotesize\renewcommand{\arraystretch}{1.15}
\begin{tabular}{lccccccccccccc}
\toprule
\textbf{Année} & 13 & 14 & 15 & 16 & 17 & 18 & 19 & 20 & 21 & 22 & 23 & 24 & 25 \\
\midrule
E1 (commune)   & - & - & $\checkmark$ & $\checkmark$ & - & $\checkmark$ & - & $\checkmark$ & $\checkmark$ & $\checkmark$ & $\checkmark$ & $\checkmark$ & $\checkmark$ \\
E2 discipl. IE & - & - & - & - & - & - & - & - & - & $\checkmark$ & $\checkmark$ & $\checkmark$ & $\checkmark$ \\
E2 appl. comm. & - & - & - & - & - & - & - & - & - & - & - & $\checkmark$ & - \\
\bottomrule
\end{tabular}
\end{center}
\smallskip
\footnotesize\textbf{Sous-thèmes fréquents :}
lecture graphique marges de Bode - réglage correcteur P ($K_{C\,\lim}$) - correcteur PI (annulation $\varepsilon_s$) - correcteur numérique (relation de récurrence) - analyse stabilité / CDC

\medskip
\textbf{Ordres de grandeur (repères rapides).}
\begin{itemize}[leftmargin=1.3em]
  \item Réseau BT FR : tension phase-neutre typique \SI{230}{\volt}, tension phase-phase typique \SI{400}{\volt}, fréquence \SI{50}{\hertz}.
  \item Courants d'étude fréquents en sujet : de quelques \si{\ampere} à quelques dizaines de \si{\ampere}.
  \item En conversion d'énergie, un rendement attendu reste généralement entre 0{,}8 et 0{,}98 selon la technologie.
\end{itemize}

\medskip
\textbf{Synthèse express (à refaire sans support).}
\begin{enumerate}[leftmargin=1.4em]
  \item Poser les hypothèses du modèle et les conventions de signe.
  \item Écrire 2 à 4 relations clés avec unités.
  \item Réaliser un calcul numérique complet et vérifier l'ordre de grandeur.
  \item Justifier le résultat par un argument technique de niveau CAPET.
\end{enumerate}

\end{tcolorbox}

\CAPETENDFILE
